\documentclass[a4paper]{article}

\newcommand{\docTitle}{Interaktive Systeme}

\usepackage{datetime2}  % for dates
\usepackage[utf8]{inputenc} % UTF-8 input encoding
\usepackage[T1]{fontenc} % fixes \textquotedbl with pdflatex/OT1
\usepackage{textcomp}   % additional text symbols for typewriter/upquote
\usepackage{xcolor}     % for gray color
\usepackage{listings}
\usepackage{marginnote} % for horizontal margin notes
\usepackage{parskip}
\usepackage{amsmath,amssymb,mathtools}
\usepackage{everypage}
\usepackage{fancyhdr}
\usepackage{amsfonts}   % Mengen and natural numbers
\usepackage{enumerate}  % Custom Enumerations
\usepackage[inline]{enumitem}   % List spacing and labels
\usepackage{multicol}   % Multiple Columns
\usepackage{tabularx}

\usepackage{hyperref}   % Links
\hypersetup{
    pdftitle={\docTitle},
    pdfauthor={Moritz Köhler},
    pdfkeywords={DHBW, Hochschule},
    colorlinks=true,
    linkcolor=black,
    filecolor=magenta,
    urlcolor=cyan
}

% -------
% Images
% -------

\usepackage{import}
\usepackage{xifthen}
\usepackage{pdfpages}
\usepackage{tikz}
\usetikzlibrary{arrows.meta, positioning}

\newcommand{\incfig}[1]{%
    \def\svgwidth{\columnwidth}
    \import{./fig/}{#1.pdf_tex}
}

% ---------------------------------
% Vertical Side note (Bottom Left)
% ---------------------------------

\newcommand{\verticaldocinfo}{%
  \rotatebox{90}{\tiny\color{gray}\texttt{\DTMtoday\_\docTitle\_\jobname}}%
}

\AddEverypageHook{
  \put(-40pt,-650pt){\verticaldocinfo}%
}

% ---------------
% Math Shortcuts
% ---------------

% Number sets
\newcommand{\R}{\mathbb{R}}
\newcommand{\N}{\mathbb{N}}
\newcommand{\Z}{\mathbb{Z}}
\newcommand{\Q}{\mathbb{Q}}
\newcommand{\C}{\mathbb{C}}

% Vectors and matrices
\newcommand{\vect}[1]{\mathbf{#1}}
\newcommand{\mat}[1]{\mathbf{#1}}

% Derivatives
\newcommand{\dd}{\mathrm{d}}
\newcommand{\dpartial}[2]{\frac{\partial #1}{\partial #2}}
\newcommand{\ddt}{\frac{\dd}{\dd t}}

% Norms and absolute values
\newcommand{\norm}[1]{\left\lVert #1 \right\rVert}
\newcommand{\abs}[1]{\left\lvert #1 \right\rvert}

% Probability & Statistics
\newcommand{\E}{\mathbb{E}}
\newcommand{\Var}{\mathrm{Var}}
\newcommand{\Prob}{\mathbb{P}}

% Fields and Algebras
\newcommand{\K}{\mathbb{K}}
\newcommand{\F}{\mathbb{F}}

% Set Operations
\newcommand{\compl}[1]{#1^{\mathrm{c}}}

% Matrix and Vector Operations
\newcommand{\T}{^{\top}}
\newcommand{\inv}{^{-1}}
\newcommand{\conj}[1]{\overline{#1}}

% -------------------
% Main lecture macro 
% -------------------

\newcommand{\lecture}[1]{%
  \protect\marginnote{\protect\tiny\protect\color{gray}#1}
}

% -----------------------------
% Command for box with heading
% -----------------------------

\fboxsep2mm
\newcommand{\know}[2]{
    \noindent\fbox{
        \parbox{\textwidth-21pt}{

            \textbf{#1:}
            \vspace{0.125cm}

            #2
        }
    }
}

% -------------------
% Listings Languages
% -------------------

% Shared defaults for code listings
\lstset{
    basicstyle=\ttfamily\small,
    keywordstyle=\color{blue}\bfseries,
    commentstyle=\color{gray},
    stringstyle=\color{teal},
    numberstyle=\tiny\color{gray},
    numbers=left,
    stepnumber=1,
    numbersep=8pt,
    showstringspaces=false,
    frame=single,
    breaklines=true,
    tabsize=4,
    keepspaces=false,
    columns=flexible,
    upquote=true,
    extendedchars=false % disable extended char handling to avoid UTF-8 conflicts
}

% SQL syntax highlighting
\lstdefinelanguage{SQL}{
    morekeywords={SELECT,FROM,WHERE,INSERT,INTO,VALUES,UPDATE,SET,DELETE,CREATE,TABLE,PRIMARY,KEY,NOT,NULL,AND,OR,JOIN,LEFT,RIGHT,INNER,ON,AS,ORDER,BY,GROUP,HAVING,LIMIT},
    sensitive=false,
    morecomment=[l]{--},
    morestring=[b]'
}

% JSON syntax highlighting
\lstdefinelanguage{json}{
    morestring=[b]",
    stringstyle=\color{teal},
    comment=[l]{//},
    morecomment=[s]{/*}{*/},
        morekeywords={true,false,null},
    literate=
     *{0}{{{\color{blue}0}}}{1}
      {1}{{{\color{blue}1}}}{1}
      {2}{{{\color{blue}2}}}{1}
      {3}{{{\color{blue}3}}}{1}
      {4}{{{\color{blue}4}}}{1}
      {5}{{{\color{blue}5}}}{1}
      {6}{{{\color{blue}6}}}{1}
      {7}{{{\color{blue}7}}}{1}
      {8}{{{\color{blue}8}}}{1}
      {9}{{{\color{blue}9}}}{1}
            {:}{{{\color{black}:}}}{1}
            {,}{{{\color{black},}}}{1},
    showstringspaces=false
}

% Language specific styles
\lstdefinestyle{javastyle}{
    language=Java,
    basicstyle=\ttfamily\small,
    keywordstyle=\color{blue}\bfseries,
    commentstyle=\color{gray},
    stringstyle=\color{teal},
    numberstyle=\tiny\color{gray},
    numbers=left,
    stepnumber=1,
    numbersep=8pt,
    showstringspaces=false,
    frame=single,
    breaklines=true,
    tabsize=4,
    keepspaces=true,
    columns=fullflexible,
    upquote=true
}

\lstdefinestyle{sqlstyle}{
    language=SQL,
    basicstyle=\ttfamily\small,
    keywordstyle=\color{blue}\bfseries,
    commentstyle=\color{gray},
    stringstyle=\color{teal},
    numberstyle=\tiny\color{gray},
    numbers=left,
    stepnumber=1,
    numbersep=8pt,
    showstringspaces=false,
    frame=single,
    breaklines=true,
    tabsize=4,
    keepspaces=true,
    columns=fullflexible,
    upquote=true
}

\lstdefinestyle{pythonstyle}{
    language=Python,
    basicstyle=\ttfamily\small,
    keywordstyle=\color{blue}\bfseries,
    commentstyle=\color{gray},
    stringstyle=\color{teal},
    numberstyle=\tiny\color{gray},
    numbers=left,
    stepnumber=1,
    numbersep=8pt,
    showstringspaces=false,
    frame=single,
    breaklines=true,
    tabsize=4,
    keepspaces=true,
    columns=fullflexible,
    upquote=true
}

\lstdefinestyle{cstyle}{
    language=C,
    basicstyle=\ttfamily\small,
    keywordstyle=\color{blue}\bfseries,
    commentstyle=\color{gray},
    stringstyle=\color{teal},
    numberstyle=\tiny\color{gray},
    numbers=left,
    stepnumber=1,
    numbersep=8pt,
    showstringspaces=false,
    frame=single,
    breaklines=true,
    tabsize=4,
    keepspaces=true,
    columns=fullflexible,
    upquote=true
}

\lstdefinestyle{bashstyle}{
    language=bash,
    basicstyle=\ttfamily\small,
    keywordstyle=\color{blue}\bfseries,
    commentstyle=\color{gray},
    stringstyle=\color{teal},
    numberstyle=\tiny\color{gray},
    numbers=left,
    stepnumber=1,
    numbersep=8pt,
    showstringspaces=false,
    frame=single,
    breaklines=true,
    tabsize=4,
    keepspaces=true,
    columns=fullflexible,
    upquote=true
}

\lstdefinestyle{jsonstyle}{
    language=json,
    basicstyle=\ttfamily\small,
    keywordstyle=\color{blue}\bfseries,
    commentstyle=\color{gray},
    stringstyle=\color{teal},
    numberstyle=\tiny\color{gray},
    numbers=left,
    stepnumber=1,
    numbersep=8pt,
    showstringspaces=false,
    frame=single,
    breaklines=true,
    tabsize=4,
    keepspaces=true,
    columns=fullflexible,
    upquote=true
}

% Default listing style
\lstset{language={}}

\title{\docTitle}
\author{Moritz}
\date{\today}

\begin{document}
\maketitle
\tableofcontents

\lecture{2026-10-01}

\section{Information}

\subsection{Menschliche Wahrnehmung}

\subsubsection{Sehen}

Der Mensch sieht Farben und Muster nur im Zentrum der Retina (des Auges)

Wir können am äußeren Rand fast nur Schwarz Weiß und Bewegungen wahrnehmen.

Blinkende Werbung ist meist immer am Rand, da diese dort am meisten Stört (Aufmerksamkeit hat)

Wir sehen auch nur in einem schmalen Kegel Scharf.

Etwas was mehr als 5 pro Sekunde blinkt, sehen wir as Punkt, nicht als blinkender Punkt.

Wir gehen durch eine Bilde in Sprüngen (Sakkade 50ms) und Fixationen. Die Reihenfolge des Bildes ist spannend und kann beeinflusst werden.

\begin{itemize}
    \item Größen und Tiefenwahrnehmung
    \item Helligkeitswahrnehmung (Luminanz und Kontrast)
    \item Farbwahrnehmung (160 Farbtöne zusätzliche weißanteile und Teilfarben insgesamt 7 Mio.)
\end{itemize}

Das Gehirn kompensiert gerne Wahrnehmungsfehler. Dies wird für optische Täuschungen genutzt und kann je nach Kontext auch zu anderen Ergebnissen führen.

\subsubsection{Höhren}

Wir können Geräusche Identifizieren (Auto, Kuh, ...), Entfernungen schätzen und die Richtung schätzen.

Wir hören Schall als Tön.

Unser \href{https://de.wikipedia.org/wiki/Lautst%C3%A4rke#/media/Datei:Akustik_db2phon.jpg}{Lautstärkeempfinden} ist dB abhängig aber nicht linear. Dies zeigt, das bei 3 kHz unseren Hören besonders empfindlich ist. Evolution hat dafür gesorgt, dass Baby Geschrei bei 3 kHz liegt.

Veränderungen von Tönen fallen deutlich auf. Angehen oder Ausschalten von z.B. der Klimaanlage.

Sprache wird deutlicher erkannt und Umgebungsgeräusche werden unterdrückt.

\subsubsection{Fühlen}

Durch die Haut: Hitze, Kälte, intensiver Druck, Schmerz.

Fühlen ist für Sehbehinderte oftmals der wichtigste Sinn.

\subsection{Schmecken und Riechen}

Hat noch keine Bedeutung für Nutzer Interfaces.

\subsection{Gestalttheorie}

Es gibt Gestaltgesetze (Gute Beispiele auf den Folien).

\begin{enumerate}
    \item Gesetz der Nähe, Unterschiede durch Ordnung/ Entfernungen - Keep the interface Simple
    \item Gesetz der Gleichheit (Form, Farbe, Größe, Orientierung)
    \item Gesetz der Geschlossenheit - Zuordnung durch Boxen oder Abstände
    \item Gesetz der Kontinuität - Die Natur machte kein Sprünge unsere Wahrnehmung auch nicht.
    \item Gesetz der Erfahrung - Das Gehirn ist Faul!
    \item Gesetz des gemeinsamen Schicksals/ Gesetz der gemeinsamen Bewegung
    \item Gesetz der Symmetrie
    \item Gesetz der Prägnanz - Einfache Strukturen werden stärker wahrgenommen als Komplizierte. Komplex wird auch einfacher im Gehirn abgespeichert.
\end{enumerate}

\subsection{Menschliche Informationsspeicherung}

\lecture{2026-10-08}

Der Mensch hat mehrere Level an Speicher:

\begin{itemize}
    \item Sensorische Speicher (Ultrakurzzeitgedächtnis)
    \item Arbeitsgedächtnis (Kurzzeitgedächtnis)
    \item Langzeitgedächtnis
\end{itemize}

Es gibt auch das Modell von Card, Moran, Newell (1983). Dies Teilt sich in drei Prozesse (gesamt 250ms):

\begin{enumerate}
    \item Wahrnehmungsprozessor (50-200ms)
    \item Kognitiver Prozessor (Unterschiedlich: 25-170ms)
    \item Motorische Prozessor (alle 30-100ms ein Neuer Bewegungsimpuls)
\end{enumerate}

Chunking und Bündelung kann genutzt werden um schneller Information Wahrzunehmen (z.B. IBAN)

\subsection{Hicks' Law und Fitt's Law}

Hick-Hyman-Gesetz:

$$t = b \log_2(n+1)$$

\begin{itemize}
    \item t: benötigte zeigt
    \item n: Anzahl der ablenkenden Elemente
    \item b: empirische Konstante (Abhängig vom Nutzer, Darstellung, lokale Dichte)
\end{itemize}

Fitts' Law

$$t = a + b*\log_2(\frac{d}{s}+1)$$

\begin{itemize}
    \item t: benötigte zeigt
    \item d: Anfangsdistanz
    \item s: Größe des Zielobjektes
    \item a: Zeit, die zum Aufnahmen des Geräts benötigt wird (Abhängig von Gerät, Alter des Nutzers)
    \item b: Geschwindigkeit des Eingabegeräte (Vgl. S. 61)
\end{itemize}

\end{document}